\documentclass[a4paper,11pt]{article}

\usepackage[top=1.5cm,bottom=1.5cm,left=2cm,right=2cm]{geometry}
\usepackage{fontspec}
\setmainfont{TeX Gyre Pagella}

\newfontfamily{\symbolfont}{Symbola}[Scale=1.0]
\newcommand{\g}[1]{{\symbolfont #1}}

\usepackage{unicode-math}
\setmathfont{TeX Gyre Pagella Math}

\usepackage{microtype}
\usepackage{array}

\pagestyle{empty}
\setlength{\parindent}{0em}
\setlength{\parskip}{0em}

\newcommand{\gAgent}{\g{▲}}
\newcommand{\gDesignar}{\g{(▲)}}
\newcommand{\gForm}{\g{●}}
\newcommand{\gGrammatikk}{\g{⬢}}
\newcommand{\gFormrom}{\g{□}}
\newcommand{\gKlasse}{\g{\{●\}}}
\newcommand{\gLyskjegle}{\g{◇}}
\newcommand{\gKonfig}{\g{⁂}}
\newcommand{\gObjekt}{\g{•}}
\newcommand{\gOmhug}{\g{☉}}
\newcommand{\gPartisjon}{\g{◐}}
\newcommand{\gTrykk}{\g{↘}}
\newcommand{\gStilart}{\g{(•••)}}
\newcommand{\gSubstrat}{\g{▽}}
\newcommand{\gLandskap}{\g{∿}}
\newcommand{\gTransform}{\g{⬟}}
\newcommand{\gSirkel}{\g{◯}}

% Seksjonsetikett inline
\newcommand{\sek}[1]{\textsc{\footnotesize #1}}

\begin{document}

% ==================== TITTEL ====================
\begin{center}
{\huge\scshape Formskrift}\\[0.2em]
{\itshape Alfabet, kompositorreglar og kjerne-likninga}
\end{center}

\vspace{0.5em}
\begin{center}\rule{\textwidth}{0.4pt}\end{center}
\vspace{0.6em}

% ==================== ALFABET ====================
\begin{center}\sek{Alfabet}\end{center}
\vspace{0.2em}

\begin{center}
\renewcommand{\arraystretch}{1.35}
\begin{tabular}{@{}c@{\hspace{0.5em}}l@{\hspace{1.8em}}c@{\hspace{0.5em}}l@{\hspace{1.8em}}c@{\hspace{0.5em}}l@{\hspace{1.8em}}c@{\hspace{0.5em}}l@{}}
\gAgent       & agent           & \gForm       & form            & \gFormrom     & formrom         & \gLyskjegle   & lyskjegle\\
\gDesignar    & designar        & \gObjekt     & objekt          & \gGrammatikk  & grammatikk      & \gOmhug       & omhug\\
\gSubstrat    & substrat        & \gKonfig     & konfigurasjon   & \gLandskap    & landskap        & \gPartisjon   & partisjon\\
\gTransform   & transformasjon  & \gKlasse     & klasse          & \gTrykk       & trykk           & \gStilart     & stilart\\
\end{tabular}
\end{center}

\vspace{0.8em}

% ==================== OPERATORAR ====================
\begin{center}\sek{Operatorar}\end{center}
\vspace{0.2em}

\begin{center}
\renewcommand{\arraystretch}{1.3}
\begin{tabular}{@{}c@{\hspace{0.4em}}l@{\hspace{1em}}c@{\hspace{0.4em}}l@{\hspace{1em}}c@{\hspace{0.4em}}l@{\hspace{1em}}c@{\hspace{0.4em}}l@{\hspace{1em}}c@{\hspace{0.4em}}l@{\hspace{1em}}c@{\hspace{0.4em}}l@{}}
$\in$             & element       & $\subset$ & delmengd    & $\cap$ & skjering & $=$ & likskap    & $\equiv$ & identitet  & $\rightarrow$ & implikasjon\\
$\leftrightarrow$ & biimplikasjon & $\triangleright$ & applikasjon & $\nabla$ & gradient & $\sum$ & sum & $|\cdot|$ & storleik & $\mid$ & vilkår\\
\end{tabular}
\end{center}

\vspace{0.8em}

% ==================== LOGIKK ====================
\begin{center}\sek{Logikk}\end{center}
\vspace{0.2em}

\begin{center}
\renewcommand{\arraystretch}{1.3}
\begin{tabular}{@{}c@{\hspace{0.5em}}l@{\hspace{1.3em}}c@{\hspace{0.5em}}l@{\hspace{1.3em}}c@{\hspace{0.5em}}l@{\hspace{1.3em}}c@{\hspace{0.5em}}l@{\hspace{1.3em}}c@{\hspace{0.5em}}l@{}}
$\forall$     & allkvantor & $\exists$ & eksistens & $\nexists$ & negert eksistens & $\wedge$ & konjunksjon & $\vee$ & disjunksjon\\
$\neg$        & negasjon   & $\therefore$ & slutning & $\because$ & premiss & $\cup$ & union & $\varnothing$ & tom mengd\\
\end{tabular}
\end{center}

\vspace{0.8em}

% ==================== KOMPOSITORREGLAR ====================
\begin{center}\sek{Kompositorreglar}\end{center}
\vspace{0.2em}

\begin{center}
\renewcommand{\arraystretch}{1.4}
\begin{tabular}{@{}r@{\hspace{0.4em}}l@{\hspace{1.5em}}r@{\hspace{0.4em}}l@{\hspace{1.5em}}r@{\hspace{0.4em}}l@{\hspace{1.5em}}r@{\hspace{0.4em}}l@{}}
\textbf{R1} & \gAgent\ $+$ \gSirkel\ $=$ \gDesignar  & \textbf{R3} & \gForm\ $+$ \gSirkel\ $=$ \gPartisjon  & \textbf{R5} & $\{\cdot\}$ $+$ \gForm\ $=$ \gKlasse    & \textbf{R7} & $\nabla$\gLandskap\ $=$ \gTrykk\\
\textbf{R2} & \gSirkel\ $+$ \gObjekt\ $=$ \gOmhug    & \textbf{R4} & \gObjekt\gObjekt\gObjekt\ $=$ \gKonfig & \textbf{R6} & $(\cdot)$ $+$ \gObjekt\gObjekt\gObjekt\ $=$ \gStilart & \textbf{R8} & \gAgent\ $+$ \gSubstrat\ $\therefore$ \gLyskjegle\\
\end{tabular}
\end{center}

\vspace{1em}
\begin{center}\rule{0.85\textwidth}{0.2pt}\end{center}
\vspace{0.8em}

% ==================== KJERNE-LIKNINGA ====================
\begin{center}\sek{Kjerne-likninga}\end{center}

\vspace{0.6em}

\begin{center}
{\Huge \gForm\ $=$ \gGrammatikk($\nabla$\gLandskap) $\cap$ \gLyskjegle}
\end{center}

\vspace{0.6em}

\begin{center}
\begin{minipage}{0.75\textwidth}
\centering\itshape\small
Forma er grammatikkens respons på landskapets gradient, projisert på agenten si lyskjegle.
\end{minipage}
\end{center}

\vspace{0.8em}
\begin{center}\rule{\textwidth}{0.4pt}\end{center}
\vspace{0.3em}

\begin{center}
\footnotesize
Seksten glyfar, tolv operatorar, ti logikk-symbol, åtte kompositorreglar. Eitt alfabet for Formlæra.
\end{center}

\end{document}
